\documentclass[11pt]{article}

% ---------------------------------------------------------------------------
% AI Applications in Research: The Case of Claude Code
% A conceptual summary compiled as a scheduled routine.
% Compiles with pdflatex, xelatex, or lualatex (only standard CTAN packages).
% ---------------------------------------------------------------------------

\usepackage[utf8]{inputenc}
\usepackage[T1]{fontenc}
\usepackage[margin=1in]{geometry}
\usepackage{microtype}
\usepackage{amsmath,amssymb}
\usepackage{enumitem}
\usepackage{booktabs}
\usepackage{xcolor}
\usepackage[colorlinks=true,linkcolor=black,urlcolor=blue,citecolor=black]{hyperref}
\usepackage{titlesec}

\titleformat{\section}{\normalfont\Large\bfseries}{\thesection}{1em}{}
\setlist[itemize]{leftmargin=1.4em,itemsep=2pt,topsep=3pt}

\title{\textbf{AI Applications in Research}\\[0.3em]
\large The Case of Claude Code and Agentic Empirical Workflows}
\author{Scheduled research summary}
\date{September 1, 2026}

\begin{document}
\maketitle

\begin{abstract}
\noindent This document summarizes the main concepts behind the use of
artificial intelligence in academic and applied research, with particular
attention to \emph{Claude Code}---a terminal-based AI coding agent---as a tool
for quantitative social science. It organizes ideas that have emerged from
practitioner accounts (notably Scott Cunningham's ``Mixtape'' series and the
accompanying \texttt{MixtapeTools} framework) into a compact conceptual map:
what the agent is, how it changes the research workflow, the verification
protocols that keep it rigorous, and the open risks that accompany its use.
It is intended as an orientation, not an endorsement or a tutorial.
\end{abstract}

\tableofcontents
\vspace{1em}

% ---------------------------------------------------------------------------
\section{Motivation}
% ---------------------------------------------------------------------------
Large language models have moved from producing isolated snippets of text or
code to operating as \emph{agents}: systems that read a repository, run
commands, inspect their own output, and iterate toward a goal. In empirical
research this shifts the unit of assistance from ``write me a function'' to
``help me carry a project from raw data to a defensible result.'' The central
claim of practitioners exploring this shift is not that AI writes the paper,
but that it changes \emph{how the human works}---the pace of iteration, the
cost of a second opinion, and the ease of replication.

The framing recurring throughout this literature is the AI as a
\textbf{thinking partner} rather than an autonomous author: the agent executes
when the next move is clear and surfaces ambiguity when it is not, so that the
\emph{dialogue itself}, not merely its output, is where the research happens.

% ---------------------------------------------------------------------------
\section{What Claude Code Is}
% ---------------------------------------------------------------------------
Claude Code is a command-line agent that operates directly inside a project
directory. Rather than a chat window detached from the work, it has tools to
read and edit files, run shell commands, execute analysis scripts, and read
back the results---closing the loop between proposing a change and observing
its effect. For a researcher this means it can:
\begin{itemize}
  \item scaffold a reproducible project structure (data, code, output, docs);
  \item translate an analysis across languages (e.g.\ \textsf{R},
        \textsf{Python}, \textsf{Stata}) and reconcile discrepancies;
  \item run, inspect, and debug estimation pipelines end to end;
  \item read long PDFs and prior literature to inform modeling choices;
  \item generate figures, tables, and presentation slides from results.
\end{itemize}

\noindent Two design features matter for research use. First, \textbf{agency
with visibility}: the agent proposes and executes concrete actions the
researcher can inspect, not opaque answers. Second, \textbf{sub-agents}: a task
can be delegated to a separate, narrowly scoped instance whose context is
isolated from the main session---a property that becomes essential for
independent verification (Section~\ref{sec:verify}).

% ---------------------------------------------------------------------------
\section{How It Changes the Research Workflow}
% ---------------------------------------------------------------------------
\begin{table}[h]
\centering
\begin{tabular}{@{}p{0.30\textwidth}p{0.62\textwidth}@{}}
\toprule
\textbf{Research stage} & \textbf{Agentic contribution} \\
\midrule
Ideation \& design & Interrogating identification strategies, restating
assumptions, proposing robustness checks. \\
Data preparation & Cleaning, merging, and documenting provenance;
flagging leakage and coding errors. \\
Estimation & Implementing estimators (e.g.\ continuous
difference-in-differences), cross-checking across languages. \\
Verification & Independent replication and structured audits by fresh
sub-agents (Section~\ref{sec:verify}). \\
Communication & Drafting slides, figures, and prose with cognitive-load and
``titles-as-assertions'' discipline. \\
\bottomrule
\end{tabular}
\caption{The agent contributes across the research lifecycle, but under human
direction and review at each stage.}
\end{table}

The workflow is explicitly \emph{iterative}: short cycles of instruct, execute,
inspect, and correct. The researcher's comparative advantage moves toward
\textbf{judgment}---posing the right question, deciding what counts as
evidence, and deciding when to stop---while mechanical execution and
boilerplate move toward the agent.

% ---------------------------------------------------------------------------
\section{Verification and Rigor}\label{sec:verify}
% ---------------------------------------------------------------------------
The most distinctive contribution of the practitioner community is a set of
\emph{verification protocols} that treat the AI's fallibility as a first-class
design constraint. The organizing insight is a \textbf{separation of powers}:
\emph{the instance that builds a pipeline cannot objectively audit it}, so
auditing is delegated to a fresh instance that never saw the original work.

\begin{description}[leftmargin=1.2em,style=nextline]
  \item[Referee 2] A multi-audit protocol (code audit, cross-language
  replication, directory structure, output automation, and econometric
  method) run \emph{after} a project is complete, in a clean context, by an
  agent that has not seen the author's reasoning. Crucially it \textbf{never
  modifies the author's code}---it writes its own replication scripts from
  scratch, so agreement is meaningful rather than circular.
  \item[Blindspot] A ``peripheral vision'' audit that names overlooked
  problems (``vices'') and missed opportunities (``virtues'')---catching what
  the author, too close to the work, cannot see.
  \item[Bibcheck] Bibliography verification by parallel, narrowly scoped agents
  that check each citation against canonical sources, guarding against
  fabricated or mis-stated references.
\end{description}

\noindent A parallel principle governs \textbf{peer review}: agents may be used
to \emph{assist} a referee---probing a paper, checking derivations, replicating
tables---without ever \emph{writing} the referee report, which remains the
reviewer's own scholarly judgment and responsibility.

% ---------------------------------------------------------------------------
\section{Communication and Presentation}
% ---------------------------------------------------------------------------
Beyond analysis, agentic tools extend to how results are communicated. The
``rhetoric of decks'' framing emphasizes beauty, deliberate management of
cognitive load, and \emph{titles as assertions} (each slide title states a
claim, not a topic). Concrete helpers include automated Beamer slide
generation and a Ti\emph{k}Z ``collision auditor'' that detects and repairs
overlapping elements in LaTeX figures---small automations that reduce the
friction between having a result and presenting it well.

% ---------------------------------------------------------------------------
\section{Risks and Open Problems}
% ---------------------------------------------------------------------------
Enthusiasm is paired with candor about failure modes. Key cautions:
\begin{itemize}
  \item \textbf{New problems for old.} Using an agent can \emph{introduce}
        errors even as it removes others, and it does not automatically save
        time---faster iteration can mean more iteration.
  \item \textbf{Automation bias.} Fluent, confident output invites
        under-scrutiny; the verification protocols above exist precisely to
        counteract this.
  \item \textbf{Circular validation.} Letting the builder grade its own work
        gives false assurance---hence the fresh-instance, no-edit rule.
  \item \textbf{Provenance and citation.} Generated references and claims must
        be checked against primary sources.
  \item \textbf{Reproducibility \& disclosure.} Because agents act in the
        environment, prompts, tool versions, and steps should be logged so
        results remain reproducible and the AI's role is transparent.
\end{itemize}

% ---------------------------------------------------------------------------
\section{Principles for Responsible Use}
% ---------------------------------------------------------------------------
A compact synthesis of the norms emerging from practice:
\begin{enumerate}[leftmargin=1.5em]
  \item Use the agent as a partner, not an oracle; keep human judgment on the
        critical path.
  \item Separate building from auditing---verify with a fresh, independent
        instance.
  \item Replicate across languages and tools; treat agreement, not a single
        run, as evidence.
  \item Check citations and data provenance against canonical sources.
  \item Log the workflow so the analysis is reproducible and the AI's
        contribution is disclosed.
  \item Reserve irreducibly human acts---the research question, the referee's
        verdict, the final claim---for the human.
\end{enumerate}

% ---------------------------------------------------------------------------
\section*{Sources and Further Reading}
% ---------------------------------------------------------------------------
\addcontentsline{toc}{section}{Sources and Further Reading}
\begin{itemize}
  \item Scott's Mixtape Substack --- \url{https://causalinf.substack.com/}
  \item Claude Code series (section) ---
        \url{https://causalinf.substack.com/s/claude-code}
  \item \texttt{scunning1975/MixtapeTools} ---
        \url{https://github.com/scunning1975/MixtapeTools}
  \item Anthropic, Claude Code --- \url{https://claude.com/claude-code}
\end{itemize}

\vfill
\noindent\rule{\linewidth}{0.4pt}\\[0.3em]
{\footnotesize This summary was compiled automatically as a scheduled routine.
Substack sources were egress-blocked at compile time; content reflects public
web-search results and the cited public repository, and should be checked
against the primary posts.}

\end{document}
